% Capítulo 2: archivo incluido desde main.tex mediante % Capítulo 2: archivo incluido desde main.tex mediante % Capítulo 2: archivo incluido desde main.tex mediante % Capítulo 2: archivo incluido desde main.tex mediante \input{chapters/02_objetivos_metodologia}.
\chapter{Objetivos y metodología}
\label{chap:objetivos-metodologia}

Este capítulo concreta qué se pretende construir, cómo se comprobará su
funcionamiento y qué evidencia será necesaria para sostener las conclusiones. La
evaluación no se plantea como una actividad posterior al desarrollo: las
condiciones de aceptación y los casos de prueba se definen antes de publicar una
versión del agente. De este modo, el TFM podrá distinguir entre una capacidad
implementada, una capacidad demostrada y una mejora pendiente.

\section{Objetivo principal}
\label{sec:objetivo-principal}

Diseñar, implementar y evaluar un agente conversacional de ciberseguridad en
Microsoft Copilot Studio, integrado en una interfaz web desarrollada con Angular,
que responda a consultas frecuentes de carácter defensivo de forma comprensible,
fundamentada en fuentes reconocidas y coherente con un conjunto explícito de
límites de seguridad.

El término \emph{evaluar} se emplea aquí en sentido estricto: se recogerán
evidencias a partir de una batería de consultas sintéticas, criterios de
aceptación predefinidos y revisión humana. No se considerará cumplido el objetivo
por el mero hecho de que el agente produzca respuestas plausibles en una
demostración aislada.

\section{Objetivos específicos}
\label{sec:objetivos-especificos}

Para alcanzar el objetivo principal se establecen los siguientes objetivos
específicos:

\begin{enumerate}
  \item Delimitar los casos de uso, actores, información permitida y situaciones
  que deben derivarse a personal especializado.
  \item Seleccionar, clasificar y documentar fuentes de conocimiento fiables de
  organismos públicos, estándares y documentación oficial de producto.
  \item Configurar el agente en Microsoft Copilot Studio con instrucciones de
  comportamiento, conocimiento acotado y respuestas de rechazo o escalado para
  solicitudes fuera de alcance.
  \item Diseñar e implementar una aplicación Angular que ofrezca una experiencia
  de consulta accesible y que integre el canal publicado sin exponer secretos ni
  debilitar los controles de autenticación.
  \item Aplicar controles básicos de seguridad a la solución, incluidos la
  minimización de datos, la gestión de configuración, la protección de la
  integración web y la trazabilidad de las fuentes.
  \item Diseñar una batería de pruebas repetible que cubra consultas normales,
  parafraseadas, ambiguas, fuera de alcance y adversariales.
  \item Analizar los resultados obtenidos, los fallos detectados y las amenazas a
  la validez de la evaluación, sin extrapolarlos a un despliegue real no
  realizado.
\end{enumerate}

\section{Requisitos del sistema}
\label{sec:requisitos}

Los requisitos se redactan de forma verificable para que puedan relacionarse con
los casos de prueba y con la evidencia de implementación. La
Tabla~\ref{tab:requisitos-funcionales} resume los requisitos funcionales
iniciales; el detalle de su implementación deberá incorporarse al capítulo
técnico junto con la referencia a la versión del artefacto evaluada.

\begin{table}[H]
\centering
\caption{Requisitos funcionales iniciales.}
\label{tab:requisitos-funcionales}
\begin{tabular}{@{}p{0.14\textwidth}p{0.61\textwidth}p{0.17\textwidth}@{}}
\hline
\textbf{ID} & \textbf{Requisito} & \textbf{Verificación prevista}\\
\hline
RF-01 & El agente debe responder consultas defensivas sobre phishing,
credenciales, dispositivos, buenas prácticas, incidentes y políticas simuladas.
& Casos por categoría.\\
RF-02 & La respuesta debe priorizar el conocimiento aprobado y, cuando la
capacidad del canal lo permita, mostrar o identificar la fuente utilizada. &
Revisión de trazabilidad.\\
RF-03 & Ante una posible incidencia, el agente debe recomendar el canal de
notificación configurado y no afirmar que ha investigado o contenido el hecho.
& Casos de derivación.\\
RF-04 & El agente debe rechazar o reconducir solicitudes ofensivas, solicitudes
de secretos y consultas ajenas a su dominio. & Casos de límites y abuso.\\
RF-05 & La interfaz Angular debe permitir iniciar una conversación, mostrar
mensajes y presentar los avisos de alcance y privacidad definidos. & Prueba
manual de interfaz.\\
RF-06 & La integración con el canal debe emplear el mecanismo de publicación y
seguridad elegido, sin incluir secretos reutilizables en código distribuido al
navegador. & Revisión de configuración.\\
\hline
\end{tabular}
\end{table}

Los requisitos no funcionales se agrupan en cinco propiedades. Primero, la
\emph{seguridad}: la aplicación no debe introducir deliberadamente una nueva vía
de exposición de secretos o datos sensibles. Segundo, la \emph{privacidad}: los
escenarios de prueba usarán información sintética y no se solicitarán
credenciales ni identificadores personales. Tercero, la \emph{usabilidad}: las
respuestas deben incluir acciones claras, advertencias proporcionadas y una
derivación inteligible. Cuarto, la \emph{trazabilidad}: ha de poder conocerse
qué fuente o versión documental justifica cada respuesta evaluada. Quinto, la
\emph{mantenibilidad}: las fuentes, instrucciones y pruebas deben poder
actualizarse sin rediseñar por completo la interfaz.

Los requisitos son una línea de base, no una garantía. En especial, la salida de
un modelo generativo depende de la configuración de la plataforma, de las fuentes
disponibles y de la formulación de la consulta. Por ello se mantendrá una lista
de cambios de configuración para explicar qué versión se ha sometido a cada
evaluación.

\section{Metodología de investigación y desarrollo}
\label{sec:metodologia}

La metodología combina desarrollo iterativo y \emph{Design Science Research}
(DSR). Este enfoque resulta apropiado porque la contribución principal es un
artefacto informático y el conocimiento se obtiene al diseñarlo y evaluarlo en
relación con un problema identificado. Hevner et al. sitúan la utilidad, el
rigor y la evaluación como principios centrales de la investigación basada en
diseño \cite{hevner2004}. Peffers et al. proponen un proceso compuesto por
identificación del problema, definición de objetivos, diseño y desarrollo,
demostración, evaluación y comunicación \cite{peffers2007}.

En este TFM dichas etapas se adaptan de la siguiente manera:

\begin{enumerate}
  \item \textbf{Identificación del problema.} Se analizan las dudas frecuentes
  de ciberseguridad y el coste de acceder a procedimientos dispersos. Se define
  el supuesto de uso organizativo simulado y los riesgos que no deben delegarse
  al agente.
  \item \textbf{Objetivos de la solución.} Se fijan los requisitos de la
  sección anterior, las categorías de consulta y los criterios de éxito. Se
  prioriza una orientación prudente y trazable frente a la amplitud temática.
  \item \textbf{Diseño y desarrollo.} Se estructura el conocimiento, se
  configuran las instrucciones y flujos del agente, y se desarrolla la interfaz
  Angular. Cada cambio relevante debe registrarse en un historial de versiones.
  \item \textbf{Demostración.} Se ejecutan conversaciones representativas desde
  el entorno configurado y desde la interfaz web. La demostración sirve para
  comprobar integración y recorrido de usuario, no para inferir calidad general.
  \item \textbf{Evaluación.} Se aplica el protocolo definido en la
  Sección~\ref{sec:protocolo-evaluacion}, se clasifican los fallos y se realiza,
  como mínimo, una iteración de mejora justificable.
  \item \textbf{Comunicación.} Se documentan arquitectura, fuentes, controles,
  resultados, limitaciones y material reproducible, distinguiendo siempre lo
  observado de lo planificado.
\end{enumerate}

La ejecución será iterativa: una evaluación puede revelar que una fuente es
ambigua, que una instrucción produce una respuesta excesivamente rígida o que una
derivación no es comprensible. La corrección debe aplicarse a la causa más
próxima que sea controlable ---por ejemplo, fuente, instrucción, flujo o
interfaz--- y después verificarse con el mismo caso de prueba. No se debe
eliminar un caso difícil para mejorar artificialmente una métrica.

\section{Gestión de fuentes de conocimiento}
\label{sec:gestion-fuentes}

La calidad del conocimiento de entrada condiciona la calidad máxima esperable de
una respuesta fundamentada. Para evitar que la base de conocimiento sea una
colección opaca de enlaces, cada elemento incorporado debe registrarse con los
siguientes atributos:

\begin{itemize}
  \item identificador, título, propietario y URL o ubicación;
  \item categoría funcional cubierta y público objetivo;
  \item fecha de publicación, fecha de última revisión conocida y fecha de
  incorporación al agente;
  \item estado de aprobación (candidato, aprobado, retirado) y responsable de la
  decisión;
  \item restricciones de acceso, idioma, licencia y datos que no deben incluirse;
  y
  \item casos de prueba asociados y versión del agente que la utiliza.
\end{itemize}

Se dará preferencia a fuentes primarias y mantenidas: marcos de NIST, informes de
ENISA, recursos de INCIBE y documentación oficial de Microsoft. Cuando una
política corporativa simulada contradiga una recomendación pública, el agente
debe expresar la prioridad de la política en su contexto y, si la contradicción
afecta a seguridad o legalidad, derivar a un responsable. Las fuentes web
públicas dependen de la indexación y de restricciones de recuperación; Microsoft
documenta que las fuentes públicas usan búsqueda basada en Bing y que existen
límites de profundidad, por lo que su configuración debe validarse con consultas
reales del conjunto de pruebas \cite{copilotPublicWeb}.

Las instrucciones también son parte de la configuración que se debe versionar.
Microsoft recomienda basarlas en herramientas y fuentes realmente configuradas,
usar nombres exactos de los recursos cuando corresponda y probar los cambios de
forma incremental \cite{copilotInstructions}. Tales recomendaciones no se
interpretan como una garantía de que el agente seguirá siempre la instrucción:
se contrastarán mediante casos de prueba.

\section{Protocolo de evaluación}
\label{sec:protocolo-evaluacion}

El protocolo persigue obtener evidencia repetible sin emplear datos reales. Un
\emph{caso de prueba} contiene, como mínimo, un identificador, categoría,
pregunta, contexto sintético, respuesta o hechos esperados, fuente de referencia,
criterio de aceptación, tipo de riesgo y resultado. Se conservará la transcripción
completa y la fecha de ejecución; si la plataforma expone la información, también
se registrará la fuente recuperada y la versión de configuración.

La batería se organizará en las cinco familias de la
Tabla~\ref{tab:familias-prueba-previstas}. Los ejemplos de la tabla son ilustrativos y no
constituyen resultados. El número definitivo de casos, la fecha de ejecución y
el identificador de versión quedan \textbf{pendientes de completar} al ejecutar
la validación real.

\begin{table}[H]
\centering
\caption{Familias de prueba previstas para la evaluación.}
\label{tab:familias-prueba-previstas}
\begin{tabular}{@{}p{0.22\textwidth}p{0.47\textwidth}p{0.23\textwidth}@{}}
\hline
\textbf{Familia} & \textbf{Propósito} & \textbf{Ejemplo de criterio}\\
\hline
Consulta directa & Verificar hechos y recomendaciones presentes en una fuente
aprobada. & Incluye acciones correctas y fuente pertinente.\\
Paráfrasis y ambigüedad & Comprobar robustez ante redacción distinta o información
insuficiente. & Mantiene el sentido o pide una aclaración mínima.\\
Incidente y escalado & Comprobar que informa, preserva el alcance y deriva sin
simular una investigación. & No solicita secretos ni promete contención.\\
Fuera de alcance & Verificar la abstención ante asuntos ajenos, jurídicos o no
fundamentados. & Declara el límite y ofrece una alternativa segura.\\
Robustez adversarial & Probar inyección de instrucciones, petición de secretos o
peticiones ofensivas. & Rechaza, no revela configuración y reconduce.\\
\hline
\end{tabular}
\end{table}

Se combinarán evaluación humana y automatización cuando la plataforma lo permita.
Copilot Studio permite crear conjuntos de prueba, comparar respuestas con
expectativas o criterios de calidad y conservar el detalle de cada caso
\cite{copilotEvaluation}. Esa capacidad es útil para repetir pruebas, pero no
elimina la revisión humana: una puntuación agregada no sustituye la comprobación
de que la fuente es pertinente, la recomendación es prudente y la derivación es
apropiada.

Para cada métrica binaria, si $n$ es el número de casos aplicables y $x_i$ vale
uno cuando el caso $i$ cumple el criterio y cero en caso contrario, se informará
la proporción de cumplimiento mediante:

\begin{equation}
M = \frac{1}{n}\sum_{i=1}^{n}x_i.
\label{eq:metrica-binaria}
\end{equation}

La ecuación~\ref{eq:metrica-binaria} se aplicará, con su denominador claramente
declarado, a las siguientes dimensiones:

\begin{itemize}
  \item \textbf{precisión}: proporción de respuestas cuyos hechos y acciones
  esenciales concuerdan con la fuente de referencia;
  \item \textbf{cobertura}: proporción de consultas de alcance que reciben una
  respuesta suficiente, no sólo parcialmente correcta;
  \item \textbf{trazabilidad}: proporción de respuestas que identifican una
  fuente pertinente y verificable cuando el caso requiere fundamentación;
  \item \textbf{consistencia}: proporción de pares de paráfrasis cuya
  recomendación esencial no se contradice;
  \item \textbf{abstención segura}: proporción de casos fuera de alcance o de
  abuso que son rechazados o derivados correctamente; y
  \item \textbf{utilidad}: valoración humana, con una rúbrica previamente
  publicada, sobre claridad, accionabilidad y proporcionalidad de la respuesta.
\end{itemize}

Las métricas se presentarán junto con el numerador, denominador y lista de casos
fallidos. No se comunicarán porcentajes sin el conjunto de pruebas, la fecha, la
configuración y los criterios que permiten interpretarlos. Antes de iniciar la
evaluación se fijarán umbrales de aceptación por familia de prueba; esos umbrales
quedan \textbf{pendientes de aprobación} por el autor y el tutor, pues dependen
del alcance final y del nivel de riesgo asumido.

\section{Plan de trabajo y entregables}
\label{sec:plan-trabajo}

La propuesta inicial distribuye el trabajo en siete semanas. La
Tabla~\ref{tab:plan-trabajo} expresa ese plan como una guía de ejecución, no como
un registro de trabajo ya completado. Las fechas de inicio, fin y entrega deben
anotarse cuando se acuerden con la dirección del TFM.

\begin{table}[H]
\centering
\caption{Plan de trabajo previsto.}
\label{tab:plan-trabajo}
\begin{tabular}{@{}p{0.14\textwidth}p{0.56\textwidth}p{0.22\textwidth}@{}}
\hline
\textbf{Fase} & \textbf{Actividad principal} & \textbf{Evidencia esperada}\\
\hline
1 & Revisión de fuentes, estado del arte y definición de alcance. & Inventario de
fuentes y requisitos.\\
2 & Diseño de casos de uso, arquitectura y experiencia de interfaz. & Diagramas,
decisiones y prototipos.\\
3 & Configuración del agente y desarrollo de la aplicación Angular. & Artefacto
ejecutable y configuración documentada.\\
4 & Integración, pruebas, corrección y evaluación controlada. & Conjunto de
pruebas, transcripciones y resultados.\\
5 & Redacción, revisión de evidencia y preparación de defensa. & Memoria,
anexos y demostración reproducible.\\
\hline
\end{tabular}
\end{table}

Los entregables incluyen el código fuente de la interfaz, la exportación o
documentación versionada del agente, el inventario de conocimiento, la batería de
pruebas, las evidencias de ejecución que no contengan información sensible y la
memoria. Las claves, secretos, identificadores de tenant, tokens y transcripciones
que contengan datos reales no deben formar parte del repositorio ni de los anexos
públicos.

\section{Aspectos éticos, de privacidad y amenazas a la validez}
\label{sec:etica-validez}

La experimentación utilizará consultas, nombres, dominios, mensajes y políticas
ficticias. La regla operativa será no introducir en el agente datos personales,
credenciales, direcciones de producción, evidencias de incidentes reales ni
información que pueda ser aprovechada para atacar un sistema. Si durante la
prueba una persona intenta aportar datos sensibles, la interfaz y las
instrucciones deben recordarle que no lo haga y el resultado no debe conservarse
sin revisión.

El RGPD exige valorar la naturaleza, el alcance, el contexto y la finalidad del
tratamiento para aplicar medidas acordes al riesgo \cite{gdpr}. El Reglamento de
IA añade obligaciones y plazos que dependen del rol, el uso y la clasificación
del sistema \cite{aiAct}. Por ello, una futura adopción real requeriría un
análisis específico de protección de datos, contratos, retención de
transcripciones, autenticación, roles, licencias y obligaciones regulatorias. El
prototipo académico no debe presentarse como evidencia suficiente de esa
conformidad.

Las principales amenazas a la validez son las siguientes:

\begin{itemize}
  \item \textbf{Validez de constructo:} una respuesta puede parecer clara y, sin
  embargo, omitir una acción crítica. La rúbrica debe separar estilo, exactitud,
  cobertura y seguridad.
  \item \textbf{Validez interna:} un cambio simultáneo de fuente, instrucción y
  modelo impide atribuir una mejora a una causa concreta. Cada iteración debe
  registrar sus modificaciones.
  \item \textbf{Validez externa:} los casos sintéticos no representan todos los
  contextos de una organización ni sustituyen pruebas con usuarios autorizados.
  \item \textbf{Reproducibilidad:} las plataformas y modelos gestionados pueden
  cambiar. La fecha, configuración disponible, fuentes y versiones reducen, pero
  no eliminan, esa limitación.
  \item \textbf{Sesgo del evaluador:} la revisión humana puede favorecer la
  expectativa del autor. Cuando sea viable, se utilizará una rúbrica y una
  segunda revisión independiente de una muestra de casos.
\end{itemize}

En consecuencia, las conclusiones deberán circunscribirse al artefacto, versión,
fuentes y batería efectivamente evaluados. La ausencia de fallos en un conjunto
finito de pruebas será evidencia de comportamiento en esos casos, no una prueba
de seguridad absoluta.

.
\chapter{Objetivos y metodología}
\label{chap:objetivos-metodologia}

Este capítulo concreta qué se pretende construir, cómo se comprobará su
funcionamiento y qué evidencia será necesaria para sostener las conclusiones. La
evaluación no se plantea como una actividad posterior al desarrollo: las
condiciones de aceptación y los casos de prueba se definen antes de publicar una
versión del agente. De este modo, el TFM podrá distinguir entre una capacidad
implementada, una capacidad demostrada y una mejora pendiente.

\section{Objetivo principal}
\label{sec:objetivo-principal}

Diseñar, implementar y evaluar un agente conversacional de ciberseguridad en
Microsoft Copilot Studio, integrado en una interfaz web desarrollada con Angular,
que responda a consultas frecuentes de carácter defensivo de forma comprensible,
fundamentada en fuentes reconocidas y coherente con un conjunto explícito de
límites de seguridad.

El término \emph{evaluar} se emplea aquí en sentido estricto: se recogerán
evidencias a partir de una batería de consultas sintéticas, criterios de
aceptación predefinidos y revisión humana. No se considerará cumplido el objetivo
por el mero hecho de que el agente produzca respuestas plausibles en una
demostración aislada.

\section{Objetivos específicos}
\label{sec:objetivos-especificos}

Para alcanzar el objetivo principal se establecen los siguientes objetivos
específicos:

\begin{enumerate}
  \item Delimitar los casos de uso, actores, información permitida y situaciones
  que deben derivarse a personal especializado.
  \item Seleccionar, clasificar y documentar fuentes de conocimiento fiables de
  organismos públicos, estándares y documentación oficial de producto.
  \item Configurar el agente en Microsoft Copilot Studio con instrucciones de
  comportamiento, conocimiento acotado y respuestas de rechazo o escalado para
  solicitudes fuera de alcance.
  \item Diseñar e implementar una aplicación Angular que ofrezca una experiencia
  de consulta accesible y que integre el canal publicado sin exponer secretos ni
  debilitar los controles de autenticación.
  \item Aplicar controles básicos de seguridad a la solución, incluidos la
  minimización de datos, la gestión de configuración, la protección de la
  integración web y la trazabilidad de las fuentes.
  \item Diseñar una batería de pruebas repetible que cubra consultas normales,
  parafraseadas, ambiguas, fuera de alcance y adversariales.
  \item Analizar los resultados obtenidos, los fallos detectados y las amenazas a
  la validez de la evaluación, sin extrapolarlos a un despliegue real no
  realizado.
\end{enumerate}

\section{Requisitos del sistema}
\label{sec:requisitos}

Los requisitos se redactan de forma verificable para que puedan relacionarse con
los casos de prueba y con la evidencia de implementación. La
Tabla~\ref{tab:requisitos-funcionales} resume los requisitos funcionales
iniciales; el detalle de su implementación deberá incorporarse al capítulo
técnico junto con la referencia a la versión del artefacto evaluada.

\begin{table}[H]
\centering
\caption{Requisitos funcionales iniciales.}
\label{tab:requisitos-funcionales}
\begin{tabular}{@{}p{0.14\textwidth}p{0.61\textwidth}p{0.17\textwidth}@{}}
\hline
\textbf{ID} & \textbf{Requisito} & \textbf{Verificación prevista}\\
\hline
RF-01 & El agente debe responder consultas defensivas sobre phishing,
credenciales, dispositivos, buenas prácticas, incidentes y políticas simuladas.
& Casos por categoría.\\
RF-02 & La respuesta debe priorizar el conocimiento aprobado y, cuando la
capacidad del canal lo permita, mostrar o identificar la fuente utilizada. &
Revisión de trazabilidad.\\
RF-03 & Ante una posible incidencia, el agente debe recomendar el canal de
notificación configurado y no afirmar que ha investigado o contenido el hecho.
& Casos de derivación.\\
RF-04 & El agente debe rechazar o reconducir solicitudes ofensivas, solicitudes
de secretos y consultas ajenas a su dominio. & Casos de límites y abuso.\\
RF-05 & La interfaz Angular debe permitir iniciar una conversación, mostrar
mensajes y presentar los avisos de alcance y privacidad definidos. & Prueba
manual de interfaz.\\
RF-06 & La integración con el canal debe emplear el mecanismo de publicación y
seguridad elegido, sin incluir secretos reutilizables en código distribuido al
navegador. & Revisión de configuración.\\
\hline
\end{tabular}
\end{table}

Los requisitos no funcionales se agrupan en cinco propiedades. Primero, la
\emph{seguridad}: la aplicación no debe introducir deliberadamente una nueva vía
de exposición de secretos o datos sensibles. Segundo, la \emph{privacidad}: los
escenarios de prueba usarán información sintética y no se solicitarán
credenciales ni identificadores personales. Tercero, la \emph{usabilidad}: las
respuestas deben incluir acciones claras, advertencias proporcionadas y una
derivación inteligible. Cuarto, la \emph{trazabilidad}: ha de poder conocerse
qué fuente o versión documental justifica cada respuesta evaluada. Quinto, la
\emph{mantenibilidad}: las fuentes, instrucciones y pruebas deben poder
actualizarse sin rediseñar por completo la interfaz.

Los requisitos son una línea de base, no una garantía. En especial, la salida de
un modelo generativo depende de la configuración de la plataforma, de las fuentes
disponibles y de la formulación de la consulta. Por ello se mantendrá una lista
de cambios de configuración para explicar qué versión se ha sometido a cada
evaluación.

\section{Metodología de investigación y desarrollo}
\label{sec:metodologia}

La metodología combina desarrollo iterativo y \emph{Design Science Research}
(DSR). Este enfoque resulta apropiado porque la contribución principal es un
artefacto informático y el conocimiento se obtiene al diseñarlo y evaluarlo en
relación con un problema identificado. Hevner et al. sitúan la utilidad, el
rigor y la evaluación como principios centrales de la investigación basada en
diseño \cite{hevner2004}. Peffers et al. proponen un proceso compuesto por
identificación del problema, definición de objetivos, diseño y desarrollo,
demostración, evaluación y comunicación \cite{peffers2007}.

En este TFM dichas etapas se adaptan de la siguiente manera:

\begin{enumerate}
  \item \textbf{Identificación del problema.} Se analizan las dudas frecuentes
  de ciberseguridad y el coste de acceder a procedimientos dispersos. Se define
  el supuesto de uso organizativo simulado y los riesgos que no deben delegarse
  al agente.
  \item \textbf{Objetivos de la solución.} Se fijan los requisitos de la
  sección anterior, las categorías de consulta y los criterios de éxito. Se
  prioriza una orientación prudente y trazable frente a la amplitud temática.
  \item \textbf{Diseño y desarrollo.} Se estructura el conocimiento, se
  configuran las instrucciones y flujos del agente, y se desarrolla la interfaz
  Angular. Cada cambio relevante debe registrarse en un historial de versiones.
  \item \textbf{Demostración.} Se ejecutan conversaciones representativas desde
  el entorno configurado y desde la interfaz web. La demostración sirve para
  comprobar integración y recorrido de usuario, no para inferir calidad general.
  \item \textbf{Evaluación.} Se aplica el protocolo definido en la
  Sección~\ref{sec:protocolo-evaluacion}, se clasifican los fallos y se realiza,
  como mínimo, una iteración de mejora justificable.
  \item \textbf{Comunicación.} Se documentan arquitectura, fuentes, controles,
  resultados, limitaciones y material reproducible, distinguiendo siempre lo
  observado de lo planificado.
\end{enumerate}

La ejecución será iterativa: una evaluación puede revelar que una fuente es
ambigua, que una instrucción produce una respuesta excesivamente rígida o que una
derivación no es comprensible. La corrección debe aplicarse a la causa más
próxima que sea controlable ---por ejemplo, fuente, instrucción, flujo o
interfaz--- y después verificarse con el mismo caso de prueba. No se debe
eliminar un caso difícil para mejorar artificialmente una métrica.

\section{Gestión de fuentes de conocimiento}
\label{sec:gestion-fuentes}

La calidad del conocimiento de entrada condiciona la calidad máxima esperable de
una respuesta fundamentada. Para evitar que la base de conocimiento sea una
colección opaca de enlaces, cada elemento incorporado debe registrarse con los
siguientes atributos:

\begin{itemize}
  \item identificador, título, propietario y URL o ubicación;
  \item categoría funcional cubierta y público objetivo;
  \item fecha de publicación, fecha de última revisión conocida y fecha de
  incorporación al agente;
  \item estado de aprobación (candidato, aprobado, retirado) y responsable de la
  decisión;
  \item restricciones de acceso, idioma, licencia y datos que no deben incluirse;
  y
  \item casos de prueba asociados y versión del agente que la utiliza.
\end{itemize}

Se dará preferencia a fuentes primarias y mantenidas: marcos de NIST, informes de
ENISA, recursos de INCIBE y documentación oficial de Microsoft. Cuando una
política corporativa simulada contradiga una recomendación pública, el agente
debe expresar la prioridad de la política en su contexto y, si la contradicción
afecta a seguridad o legalidad, derivar a un responsable. Las fuentes web
públicas dependen de la indexación y de restricciones de recuperación; Microsoft
documenta que las fuentes públicas usan búsqueda basada en Bing y que existen
límites de profundidad, por lo que su configuración debe validarse con consultas
reales del conjunto de pruebas \cite{copilotPublicWeb}.

Las instrucciones también son parte de la configuración que se debe versionar.
Microsoft recomienda basarlas en herramientas y fuentes realmente configuradas,
usar nombres exactos de los recursos cuando corresponda y probar los cambios de
forma incremental \cite{copilotInstructions}. Tales recomendaciones no se
interpretan como una garantía de que el agente seguirá siempre la instrucción:
se contrastarán mediante casos de prueba.

\section{Protocolo de evaluación}
\label{sec:protocolo-evaluacion}

El protocolo persigue obtener evidencia repetible sin emplear datos reales. Un
\emph{caso de prueba} contiene, como mínimo, un identificador, categoría,
pregunta, contexto sintético, respuesta o hechos esperados, fuente de referencia,
criterio de aceptación, tipo de riesgo y resultado. Se conservará la transcripción
completa y la fecha de ejecución; si la plataforma expone la información, también
se registrará la fuente recuperada y la versión de configuración.

La batería se organizará en las cinco familias de la
Tabla~\ref{tab:familias-prueba-previstas}. Los ejemplos de la tabla son ilustrativos y no
constituyen resultados. El número definitivo de casos, la fecha de ejecución y
el identificador de versión quedan \textbf{pendientes de completar} al ejecutar
la validación real.

\begin{table}[H]
\centering
\caption{Familias de prueba previstas para la evaluación.}
\label{tab:familias-prueba-previstas}
\begin{tabular}{@{}p{0.22\textwidth}p{0.47\textwidth}p{0.23\textwidth}@{}}
\hline
\textbf{Familia} & \textbf{Propósito} & \textbf{Ejemplo de criterio}\\
\hline
Consulta directa & Verificar hechos y recomendaciones presentes en una fuente
aprobada. & Incluye acciones correctas y fuente pertinente.\\
Paráfrasis y ambigüedad & Comprobar robustez ante redacción distinta o información
insuficiente. & Mantiene el sentido o pide una aclaración mínima.\\
Incidente y escalado & Comprobar que informa, preserva el alcance y deriva sin
simular una investigación. & No solicita secretos ni promete contención.\\
Fuera de alcance & Verificar la abstención ante asuntos ajenos, jurídicos o no
fundamentados. & Declara el límite y ofrece una alternativa segura.\\
Robustez adversarial & Probar inyección de instrucciones, petición de secretos o
peticiones ofensivas. & Rechaza, no revela configuración y reconduce.\\
\hline
\end{tabular}
\end{table}

Se combinarán evaluación humana y automatización cuando la plataforma lo permita.
Copilot Studio permite crear conjuntos de prueba, comparar respuestas con
expectativas o criterios de calidad y conservar el detalle de cada caso
\cite{copilotEvaluation}. Esa capacidad es útil para repetir pruebas, pero no
elimina la revisión humana: una puntuación agregada no sustituye la comprobación
de que la fuente es pertinente, la recomendación es prudente y la derivación es
apropiada.

Para cada métrica binaria, si $n$ es el número de casos aplicables y $x_i$ vale
uno cuando el caso $i$ cumple el criterio y cero en caso contrario, se informará
la proporción de cumplimiento mediante:

\begin{equation}
M = \frac{1}{n}\sum_{i=1}^{n}x_i.
\label{eq:metrica-binaria}
\end{equation}

La ecuación~\ref{eq:metrica-binaria} se aplicará, con su denominador claramente
declarado, a las siguientes dimensiones:

\begin{itemize}
  \item \textbf{precisión}: proporción de respuestas cuyos hechos y acciones
  esenciales concuerdan con la fuente de referencia;
  \item \textbf{cobertura}: proporción de consultas de alcance que reciben una
  respuesta suficiente, no sólo parcialmente correcta;
  \item \textbf{trazabilidad}: proporción de respuestas que identifican una
  fuente pertinente y verificable cuando el caso requiere fundamentación;
  \item \textbf{consistencia}: proporción de pares de paráfrasis cuya
  recomendación esencial no se contradice;
  \item \textbf{abstención segura}: proporción de casos fuera de alcance o de
  abuso que son rechazados o derivados correctamente; y
  \item \textbf{utilidad}: valoración humana, con una rúbrica previamente
  publicada, sobre claridad, accionabilidad y proporcionalidad de la respuesta.
\end{itemize}

Las métricas se presentarán junto con el numerador, denominador y lista de casos
fallidos. No se comunicarán porcentajes sin el conjunto de pruebas, la fecha, la
configuración y los criterios que permiten interpretarlos. Antes de iniciar la
evaluación se fijarán umbrales de aceptación por familia de prueba; esos umbrales
quedan \textbf{pendientes de aprobación} por el autor y el tutor, pues dependen
del alcance final y del nivel de riesgo asumido.

\section{Plan de trabajo y entregables}
\label{sec:plan-trabajo}

La propuesta inicial distribuye el trabajo en siete semanas. La
Tabla~\ref{tab:plan-trabajo} expresa ese plan como una guía de ejecución, no como
un registro de trabajo ya completado. Las fechas de inicio, fin y entrega deben
anotarse cuando se acuerden con la dirección del TFM.

\begin{table}[H]
\centering
\caption{Plan de trabajo previsto.}
\label{tab:plan-trabajo}
\begin{tabular}{@{}p{0.14\textwidth}p{0.56\textwidth}p{0.22\textwidth}@{}}
\hline
\textbf{Fase} & \textbf{Actividad principal} & \textbf{Evidencia esperada}\\
\hline
1 & Revisión de fuentes, estado del arte y definición de alcance. & Inventario de
fuentes y requisitos.\\
2 & Diseño de casos de uso, arquitectura y experiencia de interfaz. & Diagramas,
decisiones y prototipos.\\
3 & Configuración del agente y desarrollo de la aplicación Angular. & Artefacto
ejecutable y configuración documentada.\\
4 & Integración, pruebas, corrección y evaluación controlada. & Conjunto de
pruebas, transcripciones y resultados.\\
5 & Redacción, revisión de evidencia y preparación de defensa. & Memoria,
anexos y demostración reproducible.\\
\hline
\end{tabular}
\end{table}

Los entregables incluyen el código fuente de la interfaz, la exportación o
documentación versionada del agente, el inventario de conocimiento, la batería de
pruebas, las evidencias de ejecución que no contengan información sensible y la
memoria. Las claves, secretos, identificadores de tenant, tokens y transcripciones
que contengan datos reales no deben formar parte del repositorio ni de los anexos
públicos.

\section{Aspectos éticos, de privacidad y amenazas a la validez}
\label{sec:etica-validez}

La experimentación utilizará consultas, nombres, dominios, mensajes y políticas
ficticias. La regla operativa será no introducir en el agente datos personales,
credenciales, direcciones de producción, evidencias de incidentes reales ni
información que pueda ser aprovechada para atacar un sistema. Si durante la
prueba una persona intenta aportar datos sensibles, la interfaz y las
instrucciones deben recordarle que no lo haga y el resultado no debe conservarse
sin revisión.

El RGPD exige valorar la naturaleza, el alcance, el contexto y la finalidad del
tratamiento para aplicar medidas acordes al riesgo \cite{gdpr}. El Reglamento de
IA añade obligaciones y plazos que dependen del rol, el uso y la clasificación
del sistema \cite{aiAct}. Por ello, una futura adopción real requeriría un
análisis específico de protección de datos, contratos, retención de
transcripciones, autenticación, roles, licencias y obligaciones regulatorias. El
prototipo académico no debe presentarse como evidencia suficiente de esa
conformidad.

Las principales amenazas a la validez son las siguientes:

\begin{itemize}
  \item \textbf{Validez de constructo:} una respuesta puede parecer clara y, sin
  embargo, omitir una acción crítica. La rúbrica debe separar estilo, exactitud,
  cobertura y seguridad.
  \item \textbf{Validez interna:} un cambio simultáneo de fuente, instrucción y
  modelo impide atribuir una mejora a una causa concreta. Cada iteración debe
  registrar sus modificaciones.
  \item \textbf{Validez externa:} los casos sintéticos no representan todos los
  contextos de una organización ni sustituyen pruebas con usuarios autorizados.
  \item \textbf{Reproducibilidad:} las plataformas y modelos gestionados pueden
  cambiar. La fecha, configuración disponible, fuentes y versiones reducen, pero
  no eliminan, esa limitación.
  \item \textbf{Sesgo del evaluador:} la revisión humana puede favorecer la
  expectativa del autor. Cuando sea viable, se utilizará una rúbrica y una
  segunda revisión independiente de una muestra de casos.
\end{itemize}

En consecuencia, las conclusiones deberán circunscribirse al artefacto, versión,
fuentes y batería efectivamente evaluados. La ausencia de fallos en un conjunto
finito de pruebas será evidencia de comportamiento en esos casos, no una prueba
de seguridad absoluta.

.
\chapter{Objetivos y metodología}
\label{chap:objetivos-metodologia}

Este capítulo concreta qué se pretende construir, cómo se comprobará su
funcionamiento y qué evidencia será necesaria para sostener las conclusiones. La
evaluación no se plantea como una actividad posterior al desarrollo: las
condiciones de aceptación y los casos de prueba se definen antes de publicar una
versión del agente. De este modo, el TFM podrá distinguir entre una capacidad
implementada, una capacidad demostrada y una mejora pendiente.

\section{Objetivo principal}
\label{sec:objetivo-principal}

Diseñar, implementar y evaluar un agente conversacional de ciberseguridad en
Microsoft Copilot Studio, integrado en una interfaz web desarrollada con Angular,
que responda a consultas frecuentes de carácter defensivo de forma comprensible,
fundamentada en fuentes reconocidas y coherente con un conjunto explícito de
límites de seguridad.

El término \emph{evaluar} se emplea aquí en sentido estricto: se recogerán
evidencias a partir de una batería de consultas sintéticas, criterios de
aceptación predefinidos y revisión humana. No se considerará cumplido el objetivo
por el mero hecho de que el agente produzca respuestas plausibles en una
demostración aislada.

\section{Objetivos específicos}
\label{sec:objetivos-especificos}

Para alcanzar el objetivo principal se establecen los siguientes objetivos
específicos:

\begin{enumerate}
  \item Delimitar los casos de uso, actores, información permitida y situaciones
  que deben derivarse a personal especializado.
  \item Seleccionar, clasificar y documentar fuentes de conocimiento fiables de
  organismos públicos, estándares y documentación oficial de producto.
  \item Configurar el agente en Microsoft Copilot Studio con instrucciones de
  comportamiento, conocimiento acotado y respuestas de rechazo o escalado para
  solicitudes fuera de alcance.
  \item Diseñar e implementar una aplicación Angular que ofrezca una experiencia
  de consulta accesible y que integre el canal publicado sin exponer secretos ni
  debilitar los controles de autenticación.
  \item Aplicar controles básicos de seguridad a la solución, incluidos la
  minimización de datos, la gestión de configuración, la protección de la
  integración web y la trazabilidad de las fuentes.
  \item Diseñar una batería de pruebas repetible que cubra consultas normales,
  parafraseadas, ambiguas, fuera de alcance y adversariales.
  \item Analizar los resultados obtenidos, los fallos detectados y las amenazas a
  la validez de la evaluación, sin extrapolarlos a un despliegue real no
  realizado.
\end{enumerate}

\section{Requisitos del sistema}
\label{sec:requisitos}

Los requisitos se redactan de forma verificable para que puedan relacionarse con
los casos de prueba y con la evidencia de implementación. La
Tabla~\ref{tab:requisitos-funcionales} resume los requisitos funcionales
iniciales; el detalle de su implementación deberá incorporarse al capítulo
técnico junto con la referencia a la versión del artefacto evaluada.

\begin{table}[H]
\centering
\caption{Requisitos funcionales iniciales.}
\label{tab:requisitos-funcionales}
\begin{tabular}{@{}p{0.14\textwidth}p{0.61\textwidth}p{0.17\textwidth}@{}}
\hline
\textbf{ID} & \textbf{Requisito} & \textbf{Verificación prevista}\\
\hline
RF-01 & El agente debe responder consultas defensivas sobre phishing,
credenciales, dispositivos, buenas prácticas, incidentes y políticas simuladas.
& Casos por categoría.\\
RF-02 & La respuesta debe priorizar el conocimiento aprobado y, cuando la
capacidad del canal lo permita, mostrar o identificar la fuente utilizada. &
Revisión de trazabilidad.\\
RF-03 & Ante una posible incidencia, el agente debe recomendar el canal de
notificación configurado y no afirmar que ha investigado o contenido el hecho.
& Casos de derivación.\\
RF-04 & El agente debe rechazar o reconducir solicitudes ofensivas, solicitudes
de secretos y consultas ajenas a su dominio. & Casos de límites y abuso.\\
RF-05 & La interfaz Angular debe permitir iniciar una conversación, mostrar
mensajes y presentar los avisos de alcance y privacidad definidos. & Prueba
manual de interfaz.\\
RF-06 & La integración con el canal debe emplear el mecanismo de publicación y
seguridad elegido, sin incluir secretos reutilizables en código distribuido al
navegador. & Revisión de configuración.\\
\hline
\end{tabular}
\end{table}

Los requisitos no funcionales se agrupan en cinco propiedades. Primero, la
\emph{seguridad}: la aplicación no debe introducir deliberadamente una nueva vía
de exposición de secretos o datos sensibles. Segundo, la \emph{privacidad}: los
escenarios de prueba usarán información sintética y no se solicitarán
credenciales ni identificadores personales. Tercero, la \emph{usabilidad}: las
respuestas deben incluir acciones claras, advertencias proporcionadas y una
derivación inteligible. Cuarto, la \emph{trazabilidad}: ha de poder conocerse
qué fuente o versión documental justifica cada respuesta evaluada. Quinto, la
\emph{mantenibilidad}: las fuentes, instrucciones y pruebas deben poder
actualizarse sin rediseñar por completo la interfaz.

Los requisitos son una línea de base, no una garantía. En especial, la salida de
un modelo generativo depende de la configuración de la plataforma, de las fuentes
disponibles y de la formulación de la consulta. Por ello se mantendrá una lista
de cambios de configuración para explicar qué versión se ha sometido a cada
evaluación.

\section{Metodología de investigación y desarrollo}
\label{sec:metodologia}

La metodología combina desarrollo iterativo y \emph{Design Science Research}
(DSR). Este enfoque resulta apropiado porque la contribución principal es un
artefacto informático y el conocimiento se obtiene al diseñarlo y evaluarlo en
relación con un problema identificado. Hevner et al. sitúan la utilidad, el
rigor y la evaluación como principios centrales de la investigación basada en
diseño \cite{hevner2004}. Peffers et al. proponen un proceso compuesto por
identificación del problema, definición de objetivos, diseño y desarrollo,
demostración, evaluación y comunicación \cite{peffers2007}.

En este TFM dichas etapas se adaptan de la siguiente manera:

\begin{enumerate}
  \item \textbf{Identificación del problema.} Se analizan las dudas frecuentes
  de ciberseguridad y el coste de acceder a procedimientos dispersos. Se define
  el supuesto de uso organizativo simulado y los riesgos que no deben delegarse
  al agente.
  \item \textbf{Objetivos de la solución.} Se fijan los requisitos de la
  sección anterior, las categorías de consulta y los criterios de éxito. Se
  prioriza una orientación prudente y trazable frente a la amplitud temática.
  \item \textbf{Diseño y desarrollo.} Se estructura el conocimiento, se
  configuran las instrucciones y flujos del agente, y se desarrolla la interfaz
  Angular. Cada cambio relevante debe registrarse en un historial de versiones.
  \item \textbf{Demostración.} Se ejecutan conversaciones representativas desde
  el entorno configurado y desde la interfaz web. La demostración sirve para
  comprobar integración y recorrido de usuario, no para inferir calidad general.
  \item \textbf{Evaluación.} Se aplica el protocolo definido en la
  Sección~\ref{sec:protocolo-evaluacion}, se clasifican los fallos y se realiza,
  como mínimo, una iteración de mejora justificable.
  \item \textbf{Comunicación.} Se documentan arquitectura, fuentes, controles,
  resultados, limitaciones y material reproducible, distinguiendo siempre lo
  observado de lo planificado.
\end{enumerate}

La ejecución será iterativa: una evaluación puede revelar que una fuente es
ambigua, que una instrucción produce una respuesta excesivamente rígida o que una
derivación no es comprensible. La corrección debe aplicarse a la causa más
próxima que sea controlable ---por ejemplo, fuente, instrucción, flujo o
interfaz--- y después verificarse con el mismo caso de prueba. No se debe
eliminar un caso difícil para mejorar artificialmente una métrica.

\section{Gestión de fuentes de conocimiento}
\label{sec:gestion-fuentes}

La calidad del conocimiento de entrada condiciona la calidad máxima esperable de
una respuesta fundamentada. Para evitar que la base de conocimiento sea una
colección opaca de enlaces, cada elemento incorporado debe registrarse con los
siguientes atributos:

\begin{itemize}
  \item identificador, título, propietario y URL o ubicación;
  \item categoría funcional cubierta y público objetivo;
  \item fecha de publicación, fecha de última revisión conocida y fecha de
  incorporación al agente;
  \item estado de aprobación (candidato, aprobado, retirado) y responsable de la
  decisión;
  \item restricciones de acceso, idioma, licencia y datos que no deben incluirse;
  y
  \item casos de prueba asociados y versión del agente que la utiliza.
\end{itemize}

Se dará preferencia a fuentes primarias y mantenidas: marcos de NIST, informes de
ENISA, recursos de INCIBE y documentación oficial de Microsoft. Cuando una
política corporativa simulada contradiga una recomendación pública, el agente
debe expresar la prioridad de la política en su contexto y, si la contradicción
afecta a seguridad o legalidad, derivar a un responsable. Las fuentes web
públicas dependen de la indexación y de restricciones de recuperación; Microsoft
documenta que las fuentes públicas usan búsqueda basada en Bing y que existen
límites de profundidad, por lo que su configuración debe validarse con consultas
reales del conjunto de pruebas \cite{copilotPublicWeb}.

Las instrucciones también son parte de la configuración que se debe versionar.
Microsoft recomienda basarlas en herramientas y fuentes realmente configuradas,
usar nombres exactos de los recursos cuando corresponda y probar los cambios de
forma incremental \cite{copilotInstructions}. Tales recomendaciones no se
interpretan como una garantía de que el agente seguirá siempre la instrucción:
se contrastarán mediante casos de prueba.

\section{Protocolo de evaluación}
\label{sec:protocolo-evaluacion}

El protocolo persigue obtener evidencia repetible sin emplear datos reales. Un
\emph{caso de prueba} contiene, como mínimo, un identificador, categoría,
pregunta, contexto sintético, respuesta o hechos esperados, fuente de referencia,
criterio de aceptación, tipo de riesgo y resultado. Se conservará la transcripción
completa y la fecha de ejecución; si la plataforma expone la información, también
se registrará la fuente recuperada y la versión de configuración.

La batería se organizará en las cinco familias de la
Tabla~\ref{tab:familias-prueba-previstas}. Los ejemplos de la tabla son ilustrativos y no
constituyen resultados. El número definitivo de casos, la fecha de ejecución y
el identificador de versión quedan \textbf{pendientes de completar} al ejecutar
la validación real.

\begin{table}[H]
\centering
\caption{Familias de prueba previstas para la evaluación.}
\label{tab:familias-prueba-previstas}
\begin{tabular}{@{}p{0.22\textwidth}p{0.47\textwidth}p{0.23\textwidth}@{}}
\hline
\textbf{Familia} & \textbf{Propósito} & \textbf{Ejemplo de criterio}\\
\hline
Consulta directa & Verificar hechos y recomendaciones presentes en una fuente
aprobada. & Incluye acciones correctas y fuente pertinente.\\
Paráfrasis y ambigüedad & Comprobar robustez ante redacción distinta o información
insuficiente. & Mantiene el sentido o pide una aclaración mínima.\\
Incidente y escalado & Comprobar que informa, preserva el alcance y deriva sin
simular una investigación. & No solicita secretos ni promete contención.\\
Fuera de alcance & Verificar la abstención ante asuntos ajenos, jurídicos o no
fundamentados. & Declara el límite y ofrece una alternativa segura.\\
Robustez adversarial & Probar inyección de instrucciones, petición de secretos o
peticiones ofensivas. & Rechaza, no revela configuración y reconduce.\\
\hline
\end{tabular}
\end{table}

Se combinarán evaluación humana y automatización cuando la plataforma lo permita.
Copilot Studio permite crear conjuntos de prueba, comparar respuestas con
expectativas o criterios de calidad y conservar el detalle de cada caso
\cite{copilotEvaluation}. Esa capacidad es útil para repetir pruebas, pero no
elimina la revisión humana: una puntuación agregada no sustituye la comprobación
de que la fuente es pertinente, la recomendación es prudente y la derivación es
apropiada.

Para cada métrica binaria, si $n$ es el número de casos aplicables y $x_i$ vale
uno cuando el caso $i$ cumple el criterio y cero en caso contrario, se informará
la proporción de cumplimiento mediante:

\begin{equation}
M = \frac{1}{n}\sum_{i=1}^{n}x_i.
\label{eq:metrica-binaria}
\end{equation}

La ecuación~\ref{eq:metrica-binaria} se aplicará, con su denominador claramente
declarado, a las siguientes dimensiones:

\begin{itemize}
  \item \textbf{precisión}: proporción de respuestas cuyos hechos y acciones
  esenciales concuerdan con la fuente de referencia;
  \item \textbf{cobertura}: proporción de consultas de alcance que reciben una
  respuesta suficiente, no sólo parcialmente correcta;
  \item \textbf{trazabilidad}: proporción de respuestas que identifican una
  fuente pertinente y verificable cuando el caso requiere fundamentación;
  \item \textbf{consistencia}: proporción de pares de paráfrasis cuya
  recomendación esencial no se contradice;
  \item \textbf{abstención segura}: proporción de casos fuera de alcance o de
  abuso que son rechazados o derivados correctamente; y
  \item \textbf{utilidad}: valoración humana, con una rúbrica previamente
  publicada, sobre claridad, accionabilidad y proporcionalidad de la respuesta.
\end{itemize}

Las métricas se presentarán junto con el numerador, denominador y lista de casos
fallidos. No se comunicarán porcentajes sin el conjunto de pruebas, la fecha, la
configuración y los criterios que permiten interpretarlos. Antes de iniciar la
evaluación se fijarán umbrales de aceptación por familia de prueba; esos umbrales
quedan \textbf{pendientes de aprobación} por el autor y el tutor, pues dependen
del alcance final y del nivel de riesgo asumido.

\section{Plan de trabajo y entregables}
\label{sec:plan-trabajo}

La propuesta inicial distribuye el trabajo en siete semanas. La
Tabla~\ref{tab:plan-trabajo} expresa ese plan como una guía de ejecución, no como
un registro de trabajo ya completado. Las fechas de inicio, fin y entrega deben
anotarse cuando se acuerden con la dirección del TFM.

\begin{table}[H]
\centering
\caption{Plan de trabajo previsto.}
\label{tab:plan-trabajo}
\begin{tabular}{@{}p{0.14\textwidth}p{0.56\textwidth}p{0.22\textwidth}@{}}
\hline
\textbf{Fase} & \textbf{Actividad principal} & \textbf{Evidencia esperada}\\
\hline
1 & Revisión de fuentes, estado del arte y definición de alcance. & Inventario de
fuentes y requisitos.\\
2 & Diseño de casos de uso, arquitectura y experiencia de interfaz. & Diagramas,
decisiones y prototipos.\\
3 & Configuración del agente y desarrollo de la aplicación Angular. & Artefacto
ejecutable y configuración documentada.\\
4 & Integración, pruebas, corrección y evaluación controlada. & Conjunto de
pruebas, transcripciones y resultados.\\
5 & Redacción, revisión de evidencia y preparación de defensa. & Memoria,
anexos y demostración reproducible.\\
\hline
\end{tabular}
\end{table}

Los entregables incluyen el código fuente de la interfaz, la exportación o
documentación versionada del agente, el inventario de conocimiento, la batería de
pruebas, las evidencias de ejecución que no contengan información sensible y la
memoria. Las claves, secretos, identificadores de tenant, tokens y transcripciones
que contengan datos reales no deben formar parte del repositorio ni de los anexos
públicos.

\section{Aspectos éticos, de privacidad y amenazas a la validez}
\label{sec:etica-validez}

La experimentación utilizará consultas, nombres, dominios, mensajes y políticas
ficticias. La regla operativa será no introducir en el agente datos personales,
credenciales, direcciones de producción, evidencias de incidentes reales ni
información que pueda ser aprovechada para atacar un sistema. Si durante la
prueba una persona intenta aportar datos sensibles, la interfaz y las
instrucciones deben recordarle que no lo haga y el resultado no debe conservarse
sin revisión.

El RGPD exige valorar la naturaleza, el alcance, el contexto y la finalidad del
tratamiento para aplicar medidas acordes al riesgo \cite{gdpr}. El Reglamento de
IA añade obligaciones y plazos que dependen del rol, el uso y la clasificación
del sistema \cite{aiAct}. Por ello, una futura adopción real requeriría un
análisis específico de protección de datos, contratos, retención de
transcripciones, autenticación, roles, licencias y obligaciones regulatorias. El
prototipo académico no debe presentarse como evidencia suficiente de esa
conformidad.

Las principales amenazas a la validez son las siguientes:

\begin{itemize}
  \item \textbf{Validez de constructo:} una respuesta puede parecer clara y, sin
  embargo, omitir una acción crítica. La rúbrica debe separar estilo, exactitud,
  cobertura y seguridad.
  \item \textbf{Validez interna:} un cambio simultáneo de fuente, instrucción y
  modelo impide atribuir una mejora a una causa concreta. Cada iteración debe
  registrar sus modificaciones.
  \item \textbf{Validez externa:} los casos sintéticos no representan todos los
  contextos de una organización ni sustituyen pruebas con usuarios autorizados.
  \item \textbf{Reproducibilidad:} las plataformas y modelos gestionados pueden
  cambiar. La fecha, configuración disponible, fuentes y versiones reducen, pero
  no eliminan, esa limitación.
  \item \textbf{Sesgo del evaluador:} la revisión humana puede favorecer la
  expectativa del autor. Cuando sea viable, se utilizará una rúbrica y una
  segunda revisión independiente de una muestra de casos.
\end{itemize}

En consecuencia, las conclusiones deberán circunscribirse al artefacto, versión,
fuentes y batería efectivamente evaluados. La ausencia de fallos en un conjunto
finito de pruebas será evidencia de comportamiento en esos casos, no una prueba
de seguridad absoluta.

.
\chapter{Objetivos y metodología}
\label{chap:objetivos-metodologia}

Este capítulo concreta qué se pretende construir, cómo se comprobará su
funcionamiento y qué evidencia será necesaria para sostener las conclusiones. La
evaluación no se plantea como una actividad posterior al desarrollo: las
condiciones de aceptación y los casos de prueba se definen antes de publicar una
versión del agente. De este modo, el TFM podrá distinguir entre una capacidad
implementada, una capacidad demostrada y una mejora pendiente.

\section{Objetivo principal}
\label{sec:objetivo-principal}

Diseñar, implementar y evaluar un agente conversacional de ciberseguridad en
Microsoft Copilot Studio, integrado en una interfaz web desarrollada con Angular,
que responda a consultas frecuentes de carácter defensivo de forma comprensible,
fundamentada en fuentes reconocidas y coherente con un conjunto explícito de
límites de seguridad.

El término \emph{evaluar} se emplea aquí en sentido estricto: se recogerán
evidencias a partir de una batería de consultas sintéticas, criterios de
aceptación predefinidos y revisión humana. No se considerará cumplido el objetivo
por el mero hecho de que el agente produzca respuestas plausibles en una
demostración aislada.

\section{Objetivos específicos}
\label{sec:objetivos-especificos}

Para alcanzar el objetivo principal se establecen los siguientes objetivos
específicos:

\begin{enumerate}
  \item Delimitar los casos de uso, actores, información permitida y situaciones
  que deben derivarse a personal especializado.
  \item Seleccionar, clasificar y documentar fuentes de conocimiento fiables de
  organismos públicos, estándares y documentación oficial de producto.
  \item Configurar el agente en Microsoft Copilot Studio con instrucciones de
  comportamiento, conocimiento acotado y respuestas de rechazo o escalado para
  solicitudes fuera de alcance.
  \item Diseñar e implementar una aplicación Angular que ofrezca una experiencia
  de consulta accesible y que integre el canal publicado sin exponer secretos ni
  debilitar los controles de autenticación.
  \item Aplicar controles básicos de seguridad a la solución, incluidos la
  minimización de datos, la gestión de configuración, la protección de la
  integración web y la trazabilidad de las fuentes.
  \item Diseñar una batería de pruebas repetible que cubra consultas normales,
  parafraseadas, ambiguas, fuera de alcance y adversariales.
  \item Analizar los resultados obtenidos, los fallos detectados y las amenazas a
  la validez de la evaluación, sin extrapolarlos a un despliegue real no
  realizado.
\end{enumerate}

\section{Requisitos del sistema}
\label{sec:requisitos}

Los requisitos se redactan de forma verificable para que puedan relacionarse con
los casos de prueba y con la evidencia de implementación. La
Tabla~\ref{tab:requisitos-funcionales} resume los requisitos funcionales
iniciales; el detalle de su implementación deberá incorporarse al capítulo
técnico junto con la referencia a la versión del artefacto evaluada.

\begin{table}[H]
\centering
\caption{Requisitos funcionales iniciales.}
\label{tab:requisitos-funcionales}
\begin{tabular}{@{}p{0.14\textwidth}p{0.61\textwidth}p{0.17\textwidth}@{}}
\hline
\textbf{ID} & \textbf{Requisito} & \textbf{Verificación prevista}\\
\hline
RF-01 & El agente debe responder consultas defensivas sobre phishing,
credenciales, dispositivos, buenas prácticas, incidentes y políticas simuladas.
& Casos por categoría.\\
RF-02 & La respuesta debe priorizar el conocimiento aprobado y, cuando la
capacidad del canal lo permita, mostrar o identificar la fuente utilizada. &
Revisión de trazabilidad.\\
RF-03 & Ante una posible incidencia, el agente debe recomendar el canal de
notificación configurado y no afirmar que ha investigado o contenido el hecho.
& Casos de derivación.\\
RF-04 & El agente debe rechazar o reconducir solicitudes ofensivas, solicitudes
de secretos y consultas ajenas a su dominio. & Casos de límites y abuso.\\
RF-05 & La interfaz Angular debe permitir iniciar una conversación, mostrar
mensajes y presentar los avisos de alcance y privacidad definidos. & Prueba
manual de interfaz.\\
RF-06 & La integración con el canal debe emplear el mecanismo de publicación y
seguridad elegido, sin incluir secretos reutilizables en código distribuido al
navegador. & Revisión de configuración.\\
\hline
\end{tabular}
\end{table}

Los requisitos no funcionales se agrupan en cinco propiedades. Primero, la
\emph{seguridad}: la aplicación no debe introducir deliberadamente una nueva vía
de exposición de secretos o datos sensibles. Segundo, la \emph{privacidad}: los
escenarios de prueba usarán información sintética y no se solicitarán
credenciales ni identificadores personales. Tercero, la \emph{usabilidad}: las
respuestas deben incluir acciones claras, advertencias proporcionadas y una
derivación inteligible. Cuarto, la \emph{trazabilidad}: ha de poder conocerse
qué fuente o versión documental justifica cada respuesta evaluada. Quinto, la
\emph{mantenibilidad}: las fuentes, instrucciones y pruebas deben poder
actualizarse sin rediseñar por completo la interfaz.

Los requisitos son una línea de base, no una garantía. En especial, la salida de
un modelo generativo depende de la configuración de la plataforma, de las fuentes
disponibles y de la formulación de la consulta. Por ello se mantendrá una lista
de cambios de configuración para explicar qué versión se ha sometido a cada
evaluación.

\section{Metodología de investigación y desarrollo}
\label{sec:metodologia}

La metodología combina desarrollo iterativo y \emph{Design Science Research}
(DSR). Este enfoque resulta apropiado porque la contribución principal es un
artefacto informático y el conocimiento se obtiene al diseñarlo y evaluarlo en
relación con un problema identificado. Hevner et al. sitúan la utilidad, el
rigor y la evaluación como principios centrales de la investigación basada en
diseño \cite{hevner2004}. Peffers et al. proponen un proceso compuesto por
identificación del problema, definición de objetivos, diseño y desarrollo,
demostración, evaluación y comunicación \cite{peffers2007}.

En este TFM dichas etapas se adaptan de la siguiente manera:

\begin{enumerate}
  \item \textbf{Identificación del problema.} Se analizan las dudas frecuentes
  de ciberseguridad y el coste de acceder a procedimientos dispersos. Se define
  el supuesto de uso organizativo simulado y los riesgos que no deben delegarse
  al agente.
  \item \textbf{Objetivos de la solución.} Se fijan los requisitos de la
  sección anterior, las categorías de consulta y los criterios de éxito. Se
  prioriza una orientación prudente y trazable frente a la amplitud temática.
  \item \textbf{Diseño y desarrollo.} Se estructura el conocimiento, se
  configuran las instrucciones y flujos del agente, y se desarrolla la interfaz
  Angular. Cada cambio relevante debe registrarse en un historial de versiones.
  \item \textbf{Demostración.} Se ejecutan conversaciones representativas desde
  el entorno configurado y desde la interfaz web. La demostración sirve para
  comprobar integración y recorrido de usuario, no para inferir calidad general.
  \item \textbf{Evaluación.} Se aplica el protocolo definido en la
  Sección~\ref{sec:protocolo-evaluacion}, se clasifican los fallos y se realiza,
  como mínimo, una iteración de mejora justificable.
  \item \textbf{Comunicación.} Se documentan arquitectura, fuentes, controles,
  resultados, limitaciones y material reproducible, distinguiendo siempre lo
  observado de lo planificado.
\end{enumerate}

La ejecución será iterativa: una evaluación puede revelar que una fuente es
ambigua, que una instrucción produce una respuesta excesivamente rígida o que una
derivación no es comprensible. La corrección debe aplicarse a la causa más
próxima que sea controlable ---por ejemplo, fuente, instrucción, flujo o
interfaz--- y después verificarse con el mismo caso de prueba. No se debe
eliminar un caso difícil para mejorar artificialmente una métrica.

\section{Gestión de fuentes de conocimiento}
\label{sec:gestion-fuentes}

La calidad del conocimiento de entrada condiciona la calidad máxima esperable de
una respuesta fundamentada. Para evitar que la base de conocimiento sea una
colección opaca de enlaces, cada elemento incorporado debe registrarse con los
siguientes atributos:

\begin{itemize}
  \item identificador, título, propietario y URL o ubicación;
  \item categoría funcional cubierta y público objetivo;
  \item fecha de publicación, fecha de última revisión conocida y fecha de
  incorporación al agente;
  \item estado de aprobación (candidato, aprobado, retirado) y responsable de la
  decisión;
  \item restricciones de acceso, idioma, licencia y datos que no deben incluirse;
  y
  \item casos de prueba asociados y versión del agente que la utiliza.
\end{itemize}

Se dará preferencia a fuentes primarias y mantenidas: marcos de NIST, informes de
ENISA, recursos de INCIBE y documentación oficial de Microsoft. Cuando una
política corporativa simulada contradiga una recomendación pública, el agente
debe expresar la prioridad de la política en su contexto y, si la contradicción
afecta a seguridad o legalidad, derivar a un responsable. Las fuentes web
públicas dependen de la indexación y de restricciones de recuperación; Microsoft
documenta que las fuentes públicas usan búsqueda basada en Bing y que existen
límites de profundidad, por lo que su configuración debe validarse con consultas
reales del conjunto de pruebas \cite{copilotPublicWeb}.

Las instrucciones también son parte de la configuración que se debe versionar.
Microsoft recomienda basarlas en herramientas y fuentes realmente configuradas,
usar nombres exactos de los recursos cuando corresponda y probar los cambios de
forma incremental \cite{copilotInstructions}. Tales recomendaciones no se
interpretan como una garantía de que el agente seguirá siempre la instrucción:
se contrastarán mediante casos de prueba.

\section{Protocolo de evaluación}
\label{sec:protocolo-evaluacion}

El protocolo persigue obtener evidencia repetible sin emplear datos reales. Un
\emph{caso de prueba} contiene, como mínimo, un identificador, categoría,
pregunta, contexto sintético, respuesta o hechos esperados, fuente de referencia,
criterio de aceptación, tipo de riesgo y resultado. Se conservará la transcripción
completa y la fecha de ejecución; si la plataforma expone la información, también
se registrará la fuente recuperada y la versión de configuración.

La batería se organizará en las cinco familias de la
Tabla~\ref{tab:familias-prueba-previstas}. Los ejemplos de la tabla son ilustrativos y no
constituyen resultados. El número definitivo de casos, la fecha de ejecución y
el identificador de versión quedan \textbf{pendientes de completar} al ejecutar
la validación real.

\begin{table}[H]
\centering
\caption{Familias de prueba previstas para la evaluación.}
\label{tab:familias-prueba-previstas}
\begin{tabular}{@{}p{0.22\textwidth}p{0.47\textwidth}p{0.23\textwidth}@{}}
\hline
\textbf{Familia} & \textbf{Propósito} & \textbf{Ejemplo de criterio}\\
\hline
Consulta directa & Verificar hechos y recomendaciones presentes en una fuente
aprobada. & Incluye acciones correctas y fuente pertinente.\\
Paráfrasis y ambigüedad & Comprobar robustez ante redacción distinta o información
insuficiente. & Mantiene el sentido o pide una aclaración mínima.\\
Incidente y escalado & Comprobar que informa, preserva el alcance y deriva sin
simular una investigación. & No solicita secretos ni promete contención.\\
Fuera de alcance & Verificar la abstención ante asuntos ajenos, jurídicos o no
fundamentados. & Declara el límite y ofrece una alternativa segura.\\
Robustez adversarial & Probar inyección de instrucciones, petición de secretos o
peticiones ofensivas. & Rechaza, no revela configuración y reconduce.\\
\hline
\end{tabular}
\end{table}

Se combinarán evaluación humana y automatización cuando la plataforma lo permita.
Copilot Studio permite crear conjuntos de prueba, comparar respuestas con
expectativas o criterios de calidad y conservar el detalle de cada caso
\cite{copilotEvaluation}. Esa capacidad es útil para repetir pruebas, pero no
elimina la revisión humana: una puntuación agregada no sustituye la comprobación
de que la fuente es pertinente, la recomendación es prudente y la derivación es
apropiada.

Para cada métrica binaria, si $n$ es el número de casos aplicables y $x_i$ vale
uno cuando el caso $i$ cumple el criterio y cero en caso contrario, se informará
la proporción de cumplimiento mediante:

\begin{equation}
M = \frac{1}{n}\sum_{i=1}^{n}x_i.
\label{eq:metrica-binaria}
\end{equation}

La ecuación~\ref{eq:metrica-binaria} se aplicará, con su denominador claramente
declarado, a las siguientes dimensiones:

\begin{itemize}
  \item \textbf{precisión}: proporción de respuestas cuyos hechos y acciones
  esenciales concuerdan con la fuente de referencia;
  \item \textbf{cobertura}: proporción de consultas de alcance que reciben una
  respuesta suficiente, no sólo parcialmente correcta;
  \item \textbf{trazabilidad}: proporción de respuestas que identifican una
  fuente pertinente y verificable cuando el caso requiere fundamentación;
  \item \textbf{consistencia}: proporción de pares de paráfrasis cuya
  recomendación esencial no se contradice;
  \item \textbf{abstención segura}: proporción de casos fuera de alcance o de
  abuso que son rechazados o derivados correctamente; y
  \item \textbf{utilidad}: valoración humana, con una rúbrica previamente
  publicada, sobre claridad, accionabilidad y proporcionalidad de la respuesta.
\end{itemize}

Las métricas se presentarán junto con el numerador, denominador y lista de casos
fallidos. No se comunicarán porcentajes sin el conjunto de pruebas, la fecha, la
configuración y los criterios que permiten interpretarlos. Antes de iniciar la
evaluación se fijarán umbrales de aceptación por familia de prueba; esos umbrales
quedan \textbf{pendientes de aprobación} por el autor y el tutor, pues dependen
del alcance final y del nivel de riesgo asumido.

\section{Plan de trabajo y entregables}
\label{sec:plan-trabajo}

La propuesta inicial distribuye el trabajo en siete semanas. La
Tabla~\ref{tab:plan-trabajo} expresa ese plan como una guía de ejecución, no como
un registro de trabajo ya completado. Las fechas de inicio, fin y entrega deben
anotarse cuando se acuerden con la dirección del TFM.

\begin{table}[H]
\centering
\caption{Plan de trabajo previsto.}
\label{tab:plan-trabajo}
\begin{tabular}{@{}p{0.14\textwidth}p{0.56\textwidth}p{0.22\textwidth}@{}}
\hline
\textbf{Fase} & \textbf{Actividad principal} & \textbf{Evidencia esperada}\\
\hline
1 & Revisión de fuentes, estado del arte y definición de alcance. & Inventario de
fuentes y requisitos.\\
2 & Diseño de casos de uso, arquitectura y experiencia de interfaz. & Diagramas,
decisiones y prototipos.\\
3 & Configuración del agente y desarrollo de la aplicación Angular. & Artefacto
ejecutable y configuración documentada.\\
4 & Integración, pruebas, corrección y evaluación controlada. & Conjunto de
pruebas, transcripciones y resultados.\\
5 & Redacción, revisión de evidencia y preparación de defensa. & Memoria,
anexos y demostración reproducible.\\
\hline
\end{tabular}
\end{table}

Los entregables incluyen el código fuente de la interfaz, la exportación o
documentación versionada del agente, el inventario de conocimiento, la batería de
pruebas, las evidencias de ejecución que no contengan información sensible y la
memoria. Las claves, secretos, identificadores de tenant, tokens y transcripciones
que contengan datos reales no deben formar parte del repositorio ni de los anexos
públicos.

\section{Aspectos éticos, de privacidad y amenazas a la validez}
\label{sec:etica-validez}

La experimentación utilizará consultas, nombres, dominios, mensajes y políticas
ficticias. La regla operativa será no introducir en el agente datos personales,
credenciales, direcciones de producción, evidencias de incidentes reales ni
información que pueda ser aprovechada para atacar un sistema. Si durante la
prueba una persona intenta aportar datos sensibles, la interfaz y las
instrucciones deben recordarle que no lo haga y el resultado no debe conservarse
sin revisión.

El RGPD exige valorar la naturaleza, el alcance, el contexto y la finalidad del
tratamiento para aplicar medidas acordes al riesgo \cite{gdpr}. El Reglamento de
IA añade obligaciones y plazos que dependen del rol, el uso y la clasificación
del sistema \cite{aiAct}. Por ello, una futura adopción real requeriría un
análisis específico de protección de datos, contratos, retención de
transcripciones, autenticación, roles, licencias y obligaciones regulatorias. El
prototipo académico no debe presentarse como evidencia suficiente de esa
conformidad.

Las principales amenazas a la validez son las siguientes:

\begin{itemize}
  \item \textbf{Validez de constructo:} una respuesta puede parecer clara y, sin
  embargo, omitir una acción crítica. La rúbrica debe separar estilo, exactitud,
  cobertura y seguridad.
  \item \textbf{Validez interna:} un cambio simultáneo de fuente, instrucción y
  modelo impide atribuir una mejora a una causa concreta. Cada iteración debe
  registrar sus modificaciones.
  \item \textbf{Validez externa:} los casos sintéticos no representan todos los
  contextos de una organización ni sustituyen pruebas con usuarios autorizados.
  \item \textbf{Reproducibilidad:} las plataformas y modelos gestionados pueden
  cambiar. La fecha, configuración disponible, fuentes y versiones reducen, pero
  no eliminan, esa limitación.
  \item \textbf{Sesgo del evaluador:} la revisión humana puede favorecer la
  expectativa del autor. Cuando sea viable, se utilizará una rúbrica y una
  segunda revisión independiente de una muestra de casos.
\end{itemize}

En consecuencia, las conclusiones deberán circunscribirse al artefacto, versión,
fuentes y batería efectivamente evaluados. La ausencia de fallos en un conjunto
finito de pruebas será evidencia de comportamiento en esos casos, no una prueba
de seguridad absoluta.

